% GENERATED FILE. Edit canonical lessons, then run npm run generate:books.
% canonical-source-hash: fnv1a64:b73354bbc58c8c0f
% canonical-lessons: MW-C168-ugno, MW-C168-dhup, MW-C168-chhav, MW-C168-tulsi, MW-C168-sichno

\chapter{Sun and Shade}
\label{ch:mw-sun-and-shade}
\hlchaptermodality{168}

\begin{chapteropening}
\textbf{By the end of this chapter:} \emph{I can say when the sun rises and where the sun and the shade are.}

Complete the last lesson of chapter 168: I can say when the sun rises and where the sun and the shade are.
\end{chapteropening}

\section[\texorpdfstring{\mw{ऊगणो} (ūgṇo)}{ऊगणो (ūgṇo)}]{\mw{ऊगणो} (ūgṇo) — to rise (of the sun)}

\label{lesson:MW-C168-ugno}



\begin{quote}
Before the new one: say the Marwadi for last, then the Marwadi for at once.
\end{quote}

\subsection*{You'll want to know: \mw{ऊगणो}}
\textbf{\mw{ऊगणो}} — \emph{ūgṇo} — \textquotedblleft{}to rise (of the sun)\textquotedblright{}.

\begin{quote}
\mw{सूरज} \mw{दिनूगै} \mw{ऊगै} \mw{है।}
\end{quote}

\begin{quote}
\emph{sūraj dinūgai ūgai hai.}
\end{quote}

\begin{quote}
The sun rises early in the morning.
\end{quote}

Past (it …): \mw{वो} \mw{सूरज} \mw{ऊग्यो}, \emph{vo sūraj ūgyo}.

\begin{cousinweb}
\textbf{\mw{ऊगणो}} (\emph{ūgṇo}) is inherited from a reconstructed \emph{*udgāti}, \textquotedblleft{}rises, comes forth\textquotedblright{}; CDIAL lists Old Marwari \emph{ūgaï}, \textquotedblleft{}rises, appears\textquotedblright{} (CDIAL 1954).
\end{cousinweb}

\begin{grammarlens}[title={What you've built}]
One more word for saying what you do.
\end{grammarlens}

\subsection*{Your turn}
\begin{itemize}
\raggedright
  \item \emph{Say it:} \emph{ūgṇo}
  \item \emph{Say it:} \emph{ūgṇo}, once more
  \item \emph{Say it:} \emph{sūraj dinūgai ūgai hai.} — The sun rises early in the morning.
  \item \emph{Recall:} say the Marwadi for last, then the Marwadi for at once, then say \emph{ūgṇo} again
\end{itemize}

\begin{tcolorbox}[breakable,colback=teal!4,colframe=teal!35!black,title={Before you move on}]
What does \mw{ऊगणो} mean? (\textquotedblleft{}To rise (of the sun)\textquotedblright{}.) Say it once more.
\end{tcolorbox}
\section[\texorpdfstring{\mw{धूप} (dhūp)}{धूप (dhūp)}]{\mw{धूप} (dhūp) — sunshine}

\label{lesson:MW-C168-dhup}



\begin{quote}
Before the new one: say the Marwadi for at once, then the Marwadi for to rise.
\end{quote}

\subsection*{You'll want to know: \mw{धूप}}
\textbf{\mw{धूप}} — \emph{dhūp} — \textquotedblleft{}sunshine\textquotedblright{}.

\begin{quote}
\mw{आज} \mw{धूप} \mw{घणी} \mw{तेज} \mw{है।}
\end{quote}

\begin{quote}
\emph{āj dhūp ghaṇī tej hai.}
\end{quote}

\begin{quote}
The sun is very strong today.
\end{quote}

A feminine noun.

\begin{cousinweb}
\textbf{\mw{धूप}} (\emph{dhūp}) comes from a reconstructed \emph{*dhuppā}, \textquotedblleft{}sunshine\textquotedblright{}; CDIAL lists Hindi and Gujarati \emph{dhūp} (CDIAL 6825). CDIAL gives no Marwari form in this entry. It is feminine: \textbf{\mw{धूप} \mw{घणी} \mw{तेज} \mw{है}}.
\end{cousinweb}

\begin{grammarlens}[title={What you've built}]
One more word for saying what you do.
\end{grammarlens}

\subsection*{Your turn}
\begin{itemize}
\raggedright
  \item \emph{Say it:} \emph{dhūp}
  \item \emph{Say it:} \emph{dhūp}, once more
  \item \emph{Say it:} \emph{āj dhūp ghaṇī tej hai.} — The sun is very strong today.
  \item \emph{Recall:} say the Marwadi for at once, then the Marwadi for to rise, then say \emph{dhūp} again
\end{itemize}

\begin{tcolorbox}[breakable,colback=teal!4,colframe=teal!35!black,title={Before you move on}]
What does \mw{धूप} mean? (\textquotedblleft{}Sunshine\textquotedblright{}.) Say it once more.
\end{tcolorbox}
\section[\texorpdfstring{\mw{छांव} (chhā̃v)}{छांव (chhā̃v)}]{\mw{छांव} (chhā̃v) — shade}

\label{lesson:MW-C168-chhav}



\begin{quote}
Before the new one: say the Marwadi for to rise, then the Marwadi for sunshine.
\end{quote}

\subsection*{You'll want to know: \mw{छांव}}
\textbf{\mw{छांव}} — \emph{chhā̃v} — \textquotedblleft{}shade\textquotedblright{}.

\begin{quote}
\mw{गाय} \mw{छांव} \mw{मांय} \mw{बैठै} \mw{है।}
\end{quote}

\begin{quote}
\emph{gāy chhā̃v mā̃y baiṭhai hai.}
\end{quote}

\begin{quote}
The cow sits in the shade.
\end{quote}

A feminine noun.

\begin{cousinweb}
\textbf{\mw{छांव}} (\emph{chhā̃v}) is inherited from Sanskrit \emph{chāyā}, \textquotedblleft{}shade, shadow\textquotedblright{}; CDIAL lists Old Marwari \emph{chāvalī}, \textquotedblleft{}shade\textquotedblright{} (CDIAL 5027).
\end{cousinweb}

\begin{grammarlens}[title={What you've built}]
One more word for saying what you do.
\end{grammarlens}

\subsection*{Your turn}
\begin{itemize}
\raggedright
  \item \emph{Say it:} \emph{chhā̃v}
  \item \emph{Say it:} \emph{chhā̃v}, once more
  \item \emph{Say it:} \emph{gāy chhā̃v mā̃y baiṭhai hai.} — The cow sits in the shade.
  \item \emph{Recall:} say the Marwadi for to rise, then the Marwadi for sunshine, then say \emph{chhā̃v} again
\end{itemize}

\begin{tcolorbox}[breakable,colback=teal!4,colframe=teal!35!black,title={Before you move on}]
What does \mw{छांव} mean? (\textquotedblleft{}Shade\textquotedblright{}.) Say it once more.
\end{tcolorbox}
\section[\texorpdfstring{\mw{तुळसी} (tuḷsī)}{तुळसी (tuḷsī)}]{\mw{तुळसी} (tuḷsī) — holy basil}

\label{lesson:MW-C168-tulsi}



\begin{quote}
Before the new one: say the Marwadi for sunshine, then the Marwadi for shade.
\end{quote}

\subsection*{You'll want to know: \mw{तुळसी}}
\textbf{\mw{तुळसी}} — \emph{tuḷsī} — \textquotedblleft{}holy basil\textquotedblright{}.

\begin{quote}
\mw{म्हारी} \mw{तुळसी} \mw{हरी} \mw{अर} \mw{सुंदर} \mw{है।}
\end{quote}

\begin{quote}
\emph{mhārī tuḷsī harī ar sundar hai.}
\end{quote}

\begin{quote}
My holy basil is green and beautiful.
\end{quote}

A feminine noun.

\begin{cousinweb}
\textbf{\mw{तुळसी}} (\emph{tuḷsī}) is inherited from Sanskrit \emph{tulasī}, \textquotedblleft{}the sacred basil plant\textquotedblright{}; CDIAL lists Gujarati \emph{tuḷsī} and Hindi \emph{tulsī} (CDIAL 5885). CDIAL gives no Marwari form in this entry.
\end{cousinweb}

\begin{grammarlens}[title={What you've built}]
One more word for saying what you do.
\end{grammarlens}

\subsection*{Your turn}
\begin{itemize}
\raggedright
  \item \emph{Say it:} \emph{tuḷsī}
  \item \emph{Say it:} \emph{tuḷsī}, once more
  \item \emph{Say it:} \emph{mhārī tuḷsī harī ar sundar hai.} — My holy basil is green and beautiful.
  \item \emph{Recall:} say the Marwadi for sunshine, then the Marwadi for shade, then say \emph{tuḷsī} again
\end{itemize}

\begin{tcolorbox}[breakable,colback=teal!4,colframe=teal!35!black,title={Before you move on}]
What does \mw{तुळसी} mean? (\textquotedblleft{}Holy basil\textquotedblright{}.) Say it once more.
\end{tcolorbox}
\section[\texorpdfstring{\mw{सींचणो} (sī̃chṇo)}{सींचणो (sī̃chṇo)}]{\mw{सींचणो} (sī̃chṇo) — to water plants}

\label{lesson:MW-C168-sichno}



\begin{quote}
Before the new one: say the Marwadi for shade, then the Marwadi for holy basil.
\end{quote}

\subsection*{You'll want to know: \mw{सींचणो}}
\textbf{\mw{सींचणो}} — \emph{sī̃chṇo} — \textquotedblleft{}to water plants\textquotedblright{}.

\begin{quote}
\mw{हूं} \mw{रोज} \mw{तुळसी} \mw{सींचूं} \mw{हूं।}
\end{quote}

\begin{quote}
\emph{hū̃ roj tuḷsī sī̃chū̃ hū̃.}
\end{quote}

\begin{quote}
I water the holy basil every day.
\end{quote}

Past (I …): \mw{म्हैं} \mw{तुळसी} \mw{सींची}, \emph{mhaĩ tuḷsī sī̃chī}.

\begin{cousinweb}
\textbf{\mw{सींचणो}} (\emph{sī̃chṇo}) is inherited from Sanskrit \emph{siñcati}, \textquotedblleft{}pours out, sprinkles\textquotedblright{}; CDIAL lists Hindi \emph{sī̃cnā} and Gujarati \emph{sĩcvũ} (CDIAL 13394). CDIAL gives no Marwari form in this entry.
\end{cousinweb}

\begin{grammarlens}[title={What you've built}]
One more word for saying what you do.
\end{grammarlens}

\subsection*{Your turn}
\begin{itemize}
\raggedright
  \item \emph{Say it:} \emph{sī̃chṇo}
  \item \emph{Say it:} \emph{sī̃chṇo}, once more
  \item \emph{Say it:} \emph{hū̃ roj tuḷsī sī̃chū̃ hū̃.} — I water the holy basil every day.
  \item \emph{Recall:} say the Marwadi for shade, then the Marwadi for holy basil, then say \emph{sī̃chṇo} again
\end{itemize}

\begin{tcolorbox}[breakable,colback=teal!4,colframe=teal!35!black,title={Before you move on}]
What does \mw{सींचणो} mean? (\textquotedblleft{}To water plants\textquotedblright{}.) Say it once more.
\end{tcolorbox}
